\chapter{\texttt{auth-source}, \texttt{$\sim$/.authinfo.gpg}, and \texttt{epg-pinentry-mode}}
\label{ch:auth-source}
\minitoc\label{auth-source-.authinfo.gpg-and-epg-pinentry-mode}
\index{auth-source}\index{credentials}\index{GPG}\index{epg-pinentry-mode}\index{.authinfo.gpg}

% ── Opening signpost ──────────────────────────────────────────────
The org-mode chapters above dealt with document \emph{content};
this one deals with the \emph{credentials} that many Emacs
subsystems (email, forge, remote files) require.  This chapter explains the \texttt{auth-source} lookup
protocol, transparent GPG decryption of \texttt{\~{}/.authinfo.gpg},
and why \texttt{epg-pinentry-mode 'loopback} is required for daemon
and terminal sessions.

\section{The \texttt{auth-source} protocol}\label{the-auth-source-protocol}

\texttt{auth-source} is Emacs' credential lookup library. It searches a configurable set of backends (files, Secret Service, macOS Keychain) for entries matching \texttt{(machine HOST login USER)}. The default backend reads \texttt{$\sim$/.authinfo} or, if it exists, the GPG-encrypted \texttt{$\sim$/.authinfo.gpg}.

\texttt{(info "(auth) Top")}

\begin{minted}{elisp}
# ~/.authinfo format (one credential per line):
machine imap.example.com login alice@example.com password s3cr3t
machine smtp.example.com login alice@example.com password s3cr3t port 587
\end{minted}

Querying from Emacs Lisp:

\begin{minted}{elisp}
(auth-source-search :host "imap.example.com"
                    :user "alice@example.com"
                    :require '(:secret))
;; Returns a plist: (:host "imap.example.com" :user "…" :secret #<fn>)
;; :secret is a closure; call it with (funcall (plist-get result :secret))
\end{minted}

\section{Transparent GPG decryption}\label{transparent-gpg-decryption}

\texttt{$\sim$/.authinfo.gpg} is transparently decrypted by EasyPG (\texttt{epg}, built into Emacs) when \texttt{auth-source} reads it. No explicit decryption call is needed: \texttt{auth-source} recognises the \texttt{.gpg} extension and invokes \texttt{epa-file-find-file-hook}.

\section{\texttt{epg-pinentry-mode 'loopback}}\label{epg-pinentry-mode-loopback}

EasyPG requires a pinentry program to ask for the GPG passphrase. Three modes:

\begin{longtable}[]{@{}lll@{}}
\toprule
Mode & Passphrase prompt location & Required \\
\midrule
\texttt{default} & External pinentry GUI (pinentry-gtk, pinentry-qt, \ldots) & GTK/Qt display \\
\texttt{loopback} & Emacs minibuffer (EasyPG emulates the pinentry protocol) & Nothing external \\
\texttt{error} & Signals an error instead of prompting & (none) \\
\bottomrule
\end{longtable}

\texttt{loopback} is required when Emacs runs as a daemon (no display for an external dialog) or on a terminal Wayland session where no GTK pinentry is available:

\begin{minted}{elisp}
(use-package pinentry
  :commands magit               ; load when magit loads
  :config (pinentry-start)      ; start an in-process GPG agent
  :custom (epg-pinentry-mode 'loopback))
\end{minted}

\texttt{pinentry-start} from the \texttt{pinentry} package registers Emacs as a GPG pinentry agent via a socket, complementing \texttt{loopback} mode.

% ── Mental model summary ──────────────────────────────────────────
\section{Mental model}
The \texttt{auth-source} library searches a list of backends (files, Secret
Service, macOS Keychain) for credentials matching a given
\texttt{(machine HOST login USER)} tuple.  When the backend is
\texttt{\~{}/.authinfo.gpg}, EasyPG decrypts the file transparently so no
explicit decryption call is needed in consuming code.  Search results return
the password as a closure stored under the \texttt{:secret} key; calling it
with \texttt{funcall} yields the actual password string, keeping it out of
memory until the moment it is needed.  Setting \texttt{epg-pinentry-mode} to
\texttt{'loopback} routes the GPG passphrase prompt into the Emacs
minibuffer, which is required when no GUI pinentry program is available.
\texttt{pinentry-start} complements this by registering Emacs as a GPG agent
via a socket, so external GPG operations can also route their prompts through
Emacs.
